% Irodov, user-provided PDF, original printed pp.47--48
% (PDF pp.56--57), elasticity section problems 291--300.

\exproblem{291}{47}
玻璃薄壁管和玻璃薄壁球瓶的壁厚均为
$\Delta r=1.0\,\mathrm{mm}$，
半径均为 $r=25\,\mathrm{mm}$。
忽略外部压强，若玻璃的许用拉应力
为 $\sigma_{\max}$，分别求：
(a) 圆柱薄壁管可承受的最大内部压强；
(b) 薄壁球瓶可承受的最大内部压强。
\exsolution
假定 $\Delta r\ll r$，内部压强为 $p$。
(a) 取长为 $L$ 的圆柱管沿包含管轴的平面剖开一半，
压强在管壁法向上的合力对直径方向的投影为
$F_p=2prL$。
两条切口处管壁的拉力之和
$F_T=2\sigma\Delta r\,L$。
令二者平衡，得环向拉应力
$\sigma=pr/\Delta r$，因此
\[
\boxed{p_{\max,\mathrm{cyl}}=
\frac{\sigma_{\max}\Delta r}{r}}.
\]
(b) 对球瓶取半球，内部压强在其大圆截面上的
合力为 $p\pi r^2$，
赤道截口的薄壳拉力为
$2\pi r\Delta r\,\sigma$，故
$\sigma=pr/(2\Delta r)$，
\[
\boxed{p_{\max,\mathrm{sph}}=
\frac{2\sigma_{\max}\Delta r}{r}
=2p_{\max,\mathrm{cyl}}}.
\]
\excheck 以底本采用的玻璃拉伸强度数量级
$\sigma_{\max}\simeq5.0\times10^7\,\mathrm{Pa}$
估算，分别约 $20$ 与 $40$ 个标准大气压。
薄壁近似只适用于 $\Delta r/r=0.04\ll1$。

\exproblem{292}{47}
长度为 $\ell=1.0\,\mathrm m$ 的均匀铜杆
水平放置，绕穿过中点的竖直轴匀速旋转。
问其转速达到每秒多少转时，
因自身离心拉力在中间处发生断裂？
\exsolution
设铜杆横截面积为 $S$、密度为 $\rho$、
材料最大容许拉应力为 $\sigma_{\max}$。
在距旋转轴 $x$ 处取一质量微元
$\mathrm dm=\rho S\,\mathrm dx$，
所需向心力为
$\omega^2x\,\mathrm dm$。
中间截面必须承担整根半杆的向心力：
\[
T_{\max}=\int_0^{\ell/2}
\rho S\omega^2x\,\mathrm dx
=\frac18\rho S\omega^2\ell^2.
\]
对应的拉应力
$T_{\max}/S=\rho\omega^2\ell^2/8$。
令其等于材料强度并用
$\omega=2\pi n$，
\[
\boxed{n_{\max}=
\frac1{\pi\ell}\sqrt{\frac{2\sigma_{\max}}{\rho}}}.
\]
\excheck 结果与截面积 $S$ 无关；
以常见铜密度和抗拉强度估算约
$10^2$ 转/秒数量级，底本采用的材料数据给约
$0.8\times10^2$ 转/秒。
不同铜材的加工状态会改变 $\sigma_{\max}$，
不能只凭长度和材料名称唯一给出精确数字。

\exproblem{293}{47}
一只由铅丝制成的细圆环，
平均半径为 $r=25\,\mathrm{cm}$，
在水平面内绕穿过圆心的竖直轴匀速旋转。
已知铅的密度 $\rho$ 和许用拉应力
$\sigma_{\max}$，求圆环因旋转拉力而破裂时的转速
$n$（转/秒）。
\exsolution
设环横截面积为 $S$、线密度
$\lambda=\rho S$，
圆环以角速度 $\omega$ 转动。
取很小圆心角 $\mathrm d\varphi$ 的一段环，
其质量
$\mathrm dm=\lambda r\,\mathrm d\varphi$，
所需向心力
$\omega^2 r\,\mathrm dm=
\lambda\omega^2r^2\,\mathrm d\varphi$。
两端张力 $T$ 的向心合力一阶为
$T\,\mathrm d\varphi$，
因此 $T=\lambda\omega^2r^2$，
相应拉应力
$\sigma=T/S=\rho\omega^2r^2$。
在破裂临界时
$\rho(2\pi n)^2r^2=\sigma_{\max}$，
故
\[
\boxed{n_{\max}=\frac1{2\pi r}
\sqrt{\frac{\sigma_{\max}}{\rho}}}.
\]
\excheck 用铅的密度约
$11.3\times10^3\,\mathrm{kg/m^3}$
与底本材料强度估算约 $23$ 转/秒；
与铅丝粗细无关是均匀薄环模型的结果。

\exproblem{294}{47}
一根钢丝直径为 $d=1.0\,\mathrm{mm}$，
水平拉在相距 $\ell=2.0\,\mathrm m$ 的两固定夹点之间，
在钢丝中点悬挂质量 $m=0.25\,\mathrm{kg}$ 的重物。
求中点下降的位移 $x$，
假设弹性形变及下垂角均很小。
\exsolution
重物悬挂后，钢丝左右两半各长
$L=\sqrt{(\ell/2)^2+x^2}$。
总伸长量
\[
\Delta\ell=2L-\ell
=\ell\left[\sqrt{1+4x^2/\ell^2}-1\right]
\simeq\frac{2x^2}{\ell}.
\]
钢丝截面积 $S=\pi d^2/4$，
胡克定律给绳张力
\[
T=ES\frac{\Delta\ell}{\ell}
\simeq\frac{2ESx^2}{\ell^2}.
\]
左右两股钢丝向上的分力之和须平衡 $mg$：
\[
2T\sin\theta=mg,\qquad
\sin\theta\simeq\frac{2x}{\ell}.
\]
于是
$\frac{8ESx^3}{\ell^3}=mg$。
解得
\[
\boxed{x\simeq
\left(\frac{mg\ell^3}{8ES}\right)^{1/3}
=\left(\frac{mg\ell^3}{2\pi E d^2}\right)^{1/3}}.
\]
若钢的杨氏模量取
$E\simeq2.0\times10^{11}\,\mathrm{Pa}$，
得 $\boxed{x\simeq2.5\,\mathrm{cm}}$。
\excheck 几何伸长对下垂高度是\emph{二次}小量，
故最终位移按载荷的立方根变化，
不能套用竖直单根钢丝的线性伸长公式。

\exproblem{295}{47}
一块均匀弹性长板在光滑水平面上，
由一端面均匀分布的恒定水平合力 $F_0$
推动而整体加速。受力端面积为 $S$、
杨氏模量为 $E$。
求沿外力方向的平均纵向压缩应变。
\exsolution
设长板原长为 $\ell$、总质量为 $m$，
质量沿长度均匀分布。
整板加速度为 $a=F_0/m$。
从受力端起量坐标 $x$，
位于 $x$ 后面的剩余部分质量为
$m(\ell-x)/\ell$。
为了加速这部分质量，
截面处内部压缩力必须为
\[
N(x)=F_0\left(1-\frac{x}{\ell}\right).
\]
局部压缩应变模
$\varepsilon(x)=N(x)/(ES)$。
平均应变是沿板长的积分平均：
\[
\bar\varepsilon
=\frac1\ell\int_0^\ell
\frac{F_0}{ES}\left(1-\frac{x}{\ell}\right)\mathrm dx
=\boxed{\frac{F_0}{2ES}}.
\]
\excheck 压力从施力端 $F_0$ 线性减到另一端的零，
因此平均压缩应变是
$F_0/(ES)$ 的\emph{一半}，
而不是均匀静态轴压的数值。

\exproblem{296}{47--48}
一根长度为 $\ell$、质量为 $m$ 的均匀铜细杆
在水平面内绕一端固定的竖直轴以恒角速度
$\omega$ 转动。
求：(a) 距转轴 $r$ 处的内部拉力 $T(r)$；
(b) 整根杆因自转而产生的伸长量 $\Delta\ell$。
\exsolution
杆线密度 $\lambda=m/\ell$，
材料体密度 $\rho$、截面积 $S$，有
$m=\rho S\ell$。
距转轴 $r$ 处的横截面必须向内拉住外侧
$r\le x\le\ell$ 的所有质量微元，
各微元所需向心力为
$\omega^2x\,\lambda\,\mathrm dx$。
因此
\[
T(r)=\int_r^\ell\lambda\omega^2x\,\mathrm dx
=\frac{m\omega^2}{2\ell}(\ell^2-r^2),
\]
亦即
\[
\boxed{T(r)=\frac12m\omega^2\ell
\left(1-\frac{r^2}{\ell^2}\right)}.
\]
杆微元 $\mathrm dr$ 的伸长量为
$\mathrm d(\Delta\ell)=T(r)\mathrm dr/(ES)$，
所以
\[
\Delta\ell
=\frac{\rho\omega^2}{2E}\int_0^\ell
(\ell^2-r^2)\,\mathrm dr
=\boxed{\frac{\rho\omega^2\ell^3}{3E}}.
\]
\excheck 自由端 $r=\ell$ 处张力为零，
固定端 $r=0$ 处张力最大，
伸长与 $\omega^2$ 成正比。

\exproblem{297}{48}
一根实心铜圆柱长 $\ell=65\,\mathrm{cm}$，
置于水平面上，
其上端面均匀承受竖直向下的压缩合力
$F=1000\,\mathrm N$。
设铜的杨氏模量为 $E$、
泊松比为 $\nu$，
求圆柱体积变化量 $\Delta V$
（以 $\mathrm{mm^3}$ 表示）。
\exsolution
设圆柱端面面积为 $S$，
轴向应力 $\sigma_z=-F/S$。
各向同性材料单轴受压时，
胡克定律给
\[
\varepsilon_z=-\frac{F}{ES},\qquad
\varepsilon_x=\varepsilon_y=
-\nu\varepsilon_z=\frac{\nu F}{ES}.
\]
小变形下体积相对变化量
$\Delta V/V\simeq\varepsilon_x+\varepsilon_y+\varepsilon_z
=-(1-2\nu)F/(ES)$。
因为初始体积 $V=S\ell$，
\[
\boxed{\Delta V=-\frac{(1-2\nu)F\ell}{E}}.
\]
采用常见铜参数
$E\simeq1.2\times10^{11}\,\mathrm{Pa}$、
$\nu\simeq0.35$，
得
$\Delta V\simeq-1.63\times10^{-9}\,\mathrm{m^3}
=\boxed{-1.6\,\mathrm{mm^3}}$。
\excheck 压缩时体积缩小，因此
$\Delta V<0$（若题目只问变化量的绝对值则报告
$1.6\,\mathrm{mm^3}$）；
代数表达式中端面积 $S$ 消去。

\exproblem{298}{48}
一根密度为 $\rho$、
长度为 $\ell$ 的均匀铜杆竖直悬挂，
只由上端固定。
求：(a) 自重引起的杆总伸长量；
(b) 杆的相对体积变化量 $\Delta V/V$，
材料杨氏模量为 $E$、泊松比为 $\nu$。
\exsolution
从自由下端沿杆向上取坐标 $x$，
位于该截面下方的杆长为 $x$，
其重力导致截面拉力
$T(x)=\rho S g x$。
轴向应变
$\varepsilon_z(x)=T(x)/(ES)=\rho gx/E$。
积分杆总伸长量：
\[
\boxed{\Delta\ell=
\int_0^\ell\varepsilon_z(x)\,\mathrm dx
=\frac{\rho g\ell^2}{2E}}.
\]
单轴拉伸时横向应变为
$-\nu\varepsilon_z$，局部体积应变
$(1-2\nu)\varepsilon_z(x)$。
沿均匀杆按体积平均：
\[
\boxed{\frac{\Delta V}{V}
=(1-2\nu)\frac1\ell
\int_0^\ell\varepsilon_z(x)\,\mathrm dx
=(1-2\nu)\frac{\Delta\ell}{\ell}}.
\]
\excheck 自由端受力为零，固定端张力最大；
若将整个杆错误地视为受其总重力 $mg$ 的恒张力，
会把伸长量高估一倍。

\exproblem{299}{48}
杨氏模量为 $E$、泊松比为 $\nu$ 的各向同性均匀弹性材料
受到各向均匀压强 $p$。
求：(a) 小变形下体积相对减小量；
(b) 体积压缩系数 $\beta$ 与
$E,\nu$ 的关系；
并说明为什么稳定材料的泊松比不能大于 $1/2$。
\exsolution
静水压力状态的三个主应力相同：
$\sigma_x=\sigma_y=\sigma_z=-p$。
三维各向同性胡克定律给出每方向应变
\[
\varepsilon_x=
\frac1E[\sigma_x-\nu(\sigma_y+\sigma_z)]
=-\frac{p}{E}(1-2\nu),
\]
且 $\varepsilon_y=\varepsilon_z=\varepsilon_x$。
小变形下体积应变为三者相加：
\[
\boxed{\frac{\Delta V}{V}
=-\frac{3p}{E}(1-2\nu)}.
\]
若把体积\emph{减小的比例}定义为正，
则为 $3p(1-2\nu)/E$。
压缩系数定义
$\beta=-(1/V)\mathrm dV/\mathrm dp$，
故
\[
\boxed{\beta=\frac{3(1-2\nu)}{E}}.
\]
稳定材料受正压强时不能自发膨胀，
要求 $\beta\ge0$；
在 $E>0$ 情形这意味着
$\boxed{\nu\le1/2}$。
\excheck 对严格稳定且可压缩的普通弹性固体，
通常有 $\nu<1/2$；
$\nu=1/2$ 是理想不可压缩极限。

\exproblem{300}{48}
一根钢质均匀矩形悬臂梁，一端固定在墙中，
伸出长度为 $\ell=6.0\,\mathrm m$，
竖直厚度为 $h=10\,\mathrm{cm}$，
宽度为 $b$（原图1.74）。
梁仅因自重发生微小下挠。
求靠近固定端 $O$ 处中性层的曲率半径 $R_c$。
\exsolution
取梁宽度为 $b$、厚度为 $h$，
横截面积 $S=bh$，线重
$q=\rho S g=\rho b h g$，
其中 $\rho$ 为钢密度。
固定端截面承受右侧整根悬臂均匀重载引起的弯矩：
\[
M_O=\int_0^\ell q x\,\mathrm dx
=\frac12q\ell^2.
\]
矩形截面对水平中性轴的截面二次矩
$I_A=bh^3/12$。
在小挠度的欧拉—伯努利梁模型中，
曲率 $\kappa\simeq|M_O|/(EI_A)$，
因此
\[
R_c=\frac1\kappa
=\frac{EI_A}{M_O}
=\frac{E(bh^3/12)}{(\rho bhg)\ell^2/2}
=\boxed{\frac{Eh^2}{6\rho g\ell^2}}.
\]
采用钢的常见参数
$E\simeq2.0\times10^{11}\,\mathrm{Pa}$、
$\rho\simeq7.8\times10^3\,\mathrm{kg/m^3}$，
得
$\boxed{R_c\simeq0.12\,\mathrm{km}}$。
\excheck 梁宽 $b$ 完全约去，曲率半径不依赖宽度。
结果仅适用于小挠度、线弹性、
等截面与重载沿梁长均匀的近似。
